\documentclass{amsart}
% For dealing with references we use the comment environment
\usepackage{verbatim}
\newenvironment{reference}{\comment}{\endcomment}
%\newenvironment{reference}{}{}
\newenvironment{slogan}{\comment}{\endcomment}
\newenvironment{history}{\comment}{\endcomment}

% For commutative diagrams we use Xy-pic
\usepackage[all]{xy}

% We use 2cell for 2-commutative diagrams.
\xyoption{2cell}
\UseAllTwocells

% We use multicol for the list of chapters between chapters
\usepackage{multicol}

% This is generally recommended for better output
\usepackage{lmodern}
\usepackage[T1]{fontenc}

% For cross-file-references

% Package for hypertext links:
\usepackage{hyperref}

% For any local file, say "hello.tex" you want to link to please

% Theorem environments.
%
\theoremstyle{plain}
\newtheorem{theorem}[subsection]{Theorem}
\newtheorem{proposition}[subsection]{Proposition}
\newtheorem{lemma}[subsection]{Lemma}

\theoremstyle{definition}
\newtheorem{definition}[subsection]{Definition}
\newtheorem{example}[subsection]{Example}
\newtheorem{exercise}[subsection]{Exercise}
\newtheorem{situation}[subsection]{Situation}

\theoremstyle{remark}
\newtheorem{remark}[subsection]{Remark}
\newtheorem{remarks}[subsection]{Remarks}

\numberwithin{equation}{subsection}

% Macros
%
\def\lim{\mathop{\mathrm{lim}}\nolimits}
\def\colim{\mathop{\mathrm{colim}}\nolimits}
\def\Spec{\mathop{\mathrm{Spec}}}
\def\Hom{\mathop{\mathrm{Hom}}\nolimits}
\def\Ext{\mathop{\mathrm{Ext}}\nolimits}
\def\SheafHom{\mathop{\mathcal{H}\!\mathit{om}}\nolimits}
\def\SheafExt{\mathop{\mathcal{E}\!\mathit{xt}}\nolimits}
\def\Sch{\mathit{Sch}}
\def\Mor{\mathop{\mathrm{Mor}}\nolimits}
\def\Ob{\mathop{\mathrm{Ob}}\nolimits}
\def\Sh{\mathop{\mathit{Sh}}\nolimits}
\def\NL{\mathop{N\!L}\nolimits}
\def\CH{\mathop{\mathrm{CH}}\nolimits}
\def\proetale{{pro\text{-}\acute{e}tale}}
\def\etale{{\acute{e}tale}}
\def\QCoh{\mathit{QCoh}}
\def\Ker{\mathop{\mathrm{Ker}}}
\def\Im{\mathop{\mathrm{Im}}}
\def\Coker{\mathop{\mathrm{Coker}}}
\def\Coim{\mathop{\mathrm{Coim}}}

% Boxtimes
%
\DeclareMathSymbol{\boxtimes}{\mathbin}{AMSa}{"02}

%
% Macros for moduli stacks/spaces
%
\def\QCohstack{\mathcal{QC}\!\mathit{oh}}
\def\Cohstack{\mathcal{C}\!\mathit{oh}}
\def\Spacesstack{\mathcal{S}\!\mathit{paces}}
\def\Quotfunctor{\mathrm{Quot}}
\def\Hilbfunctor{\mathrm{Hilb}}
\def\Curvesstack{\mathcal{C}\!\mathit{urves}}
\def\Polarizedstack{\mathcal{P}\!\mathit{olarized}}
\def\Complexesstack{\mathcal{C}\!\mathit{omplexes}}
% \Pic is the operator that assigns to X its picard group, usage \Pic(X)
% \Picardstack_{X/B} denotes the Picard stack of X over B
% \Picardfunctor_{X/B} denotes the Picard functor of X over B
\def\Pic{\mathop{\mathrm{Pic}}\nolimits}
\def\Picardstack{\mathcal{P}\!\mathit{ic}}
\def\Picardfunctor{\mathrm{Pic}}
\def\Deformationcategory{\mathcal{D}\!\mathit{ef}}

\setlength{\emergencystretch}{3em}
\begin{document}
\title[Stacks Project, Tag 0GN3]{655 --- Global cohomology commutes with filtered colimits under compact generators}
\maketitle
\noindent Copyright (C) 2005--2025 Johan de Jong. Stacks Project authors. Source: \href{https://stacks.math.columbia.edu/tag/0GN3}{Tag 0GN3}.

\noindent Editorial difficulty note (not author text). Difficulty H2: The proof uses a spectral sequence over powers of a covering sheaf and verifies the generator hypotheses anew in every slice topos. Iterated equalizer refinements and simultaneous induction over topoi add substantial content beyond the basic space-level colimit theorem..

\noindent Editorial provenance note (not author text). History: excerpt from revision \texttt{a0444\allowbreak 6e57e\allowbreak c1fbc\allowbreak 252a8\allowbreak 71afc\allowbreak ec775\allowbreak 2fb28\allowbreak 07b14}, sites-cohomology.tex, lines 2629--2785. Reference presentation was adapted; original-author context and core support blocks were retained or added. The mathematical wording is unchanged. Licensed under GNU FDL 1.2 or later; no invariant sections or cover texts. See ../licenses/COPYING. Context blocks reproduce author definitions and supporting results; they are not additional counted problems.

\begin{definition}
\label{sites-definition-site}
A {\it site}\footnote{This notation differs from that of \href{https://stacks.math.columbia.edu/bibliography}{Bibliography: SGA4}, as
explained in the introduction.} is given by a category $\mathcal{C}$ and a set
$\text{Cov}(\mathcal{C})$ of families of morphisms with fixed target
$\{U_i \to U\}_{i \in I}$, called {\it coverings of $\mathcal{C}$},
satisfying the following axioms
\begin{enumerate}
\item If $V \to U$ is an isomorphism then $\{V \to U\} \in
\text{Cov}(\mathcal{C})$.
\item If $\{U_i \to U\}_{i\in I} \in \text{Cov}(\mathcal{C})$ and for each
$i$ we have $\{V_{ij} \to U_i\}_{j\in J_i} \in \text{Cov}(\mathcal{C})$, then
$\{V_{ij} \to U\}_{i \in I, j\in J_i} \in \text{Cov}(\mathcal{C})$.
\item If $\{U_i \to U\}_{i\in I}\in \text{Cov}(\mathcal{C})$
and $V \to U$ is a morphism of $\mathcal{C}$ then $U_i \times_U V$
exists for all $i$ and
$\{U_i \times_U V \to V \}_{i\in I} \in \text{Cov}(\mathcal{C})$.
\end{enumerate}
\end{definition}

\begin{definition}
\label{sites-definition-quasi-compact-topos}
An object $\mathcal{F}$ of a topos $\Sh(\mathcal{C})$ is {\it quasi-compact}
if for any surjective map $\coprod_{i \in I} \mathcal{F}_i \to \mathcal{F}$
of $\Sh(\mathcal{C})$ there exists a finite subset $J \subset I$ such
that $\coprod_{i \in J} \mathcal{F}_i \to \mathcal{F}$ is surjective.
A topos $\Sh(\mathcal{C})$ is said to be {\it quasi-compact}
if its final object $*$ is a quasi-compact object.
\end{definition}

\begin{lemma}
\label{lemma-colim-works-over-collection}
Let $\mathcal{C}$ be a site. Let $\text{Cov}_\mathcal{C}$ be the set
of coverings of $\mathcal{C}$ (see
Sites, Definition \ref{sites-definition-site}). Let
$\mathcal{B} \subset \Ob(\mathcal{C})$, and
$\text{Cov} \subset \text{Cov}_\mathcal{C}$
be subsets. Assume that
\begin{enumerate}
\item For every $\mathcal{U} \in \text{Cov}$ we have
$\mathcal{U} = \{U_i \to U\}_{i \in I}$ with $I$ finite,
$U, U_i \in \mathcal{B}$ and every
$U_{i_0} \times_U \ldots \times_U U_{i_p} \in \mathcal{B}$.
\item For every $U \in \mathcal{B}$ the coverings of $U$
occurring in $\text{Cov}$ is a cofinal system of coverings of $U$.
\end{enumerate}
Then the map
$$
\colim_i H^p(U, \mathcal{F}_i)
\longrightarrow
H^p(U, \colim_i \mathcal{F}_i)
$$
is an isomorphism for every $p \geq 0$, every $U \in \mathcal{B}$, and
every filtered diagram $\mathcal{I} \to \textit{Ab}(\mathcal{C})$.
\end{lemma}

\begin{proof}
To prove the lemma we will argue by induction on $p$.
Note that we require in (1) the coverings $\mathcal{U} \in \text{Cov}$
to be finite, so that all the elements of $\mathcal{B}$ are quasi-compact.
Hence (2) and (1) imply that any $U \in \mathcal{B}$ satisfies the hypothesis
of Sites, Lemma \href{https://stacks.math.columbia.edu/tag/0738}{lemma directed colimits sections (Tag 0738)} (4).
Thus we see that the result holds for $p = 0$.
Now we assume the lemma holds for $p$ and prove it for $p + 1$.

\medskip\noindent
Choose a filtered diagram
$\mathcal{F} : \mathcal{I} \to \textit{Ab}(\mathcal{C})$,
$i \mapsto \mathcal{F}_i$.
Since $\textit{Ab}(\mathcal{C})$ has functorial injective embeddings, see
Injectives, Theorem \href{https://stacks.math.columbia.edu/tag/01DP}{theorem sheaves injectives (Tag 01DP)},
we can find a morphism of filtered diagrams
$\mathcal{F} \to \mathcal{I}$
such that each $\mathcal{F}_i \to \mathcal{I}_i$ is an injective map of
abelian sheaves into an injective abelian sheaf. Denote $\mathcal{Q}_i$
the cokernel so that we have short exact sequences
$$
0 \to
\mathcal{F}_i \to
\mathcal{I}_i \to
\mathcal{Q}_i \to 0.
$$
Since colimits of sheaves are the sheafification of colimits on the level
of presheaves, since sheafification is exact, and since filtered
colimits of abelian groups are exact
(see Algebra, Lemma \href{https://stacks.math.columbia.edu/tag/00DB}{lemma directed colimit exact (Tag 00DB)}),
we see the sequence
$$
0 \to
\colim_i \mathcal{F}_i \to
\colim_i \mathcal{I}_i \to
\colim_i \mathcal{Q}_i \to 0.
$$
is also a short exact sequence. We claim that
$H^q(U, \colim_i \mathcal{I}_i) = 0$ for all $U \in \mathcal{B}$
and all $q \geq 1$. Accepting this claim
for the moment consider the diagram
$$
\xymatrix{
\colim_i H^p(U, \mathcal{I}_i) \ar[d] \ar[r] &
\colim_i H^p(U, \mathcal{Q}_i) \ar[d] \ar[r] &
\colim_i H^{p + 1}(U, \mathcal{F}_i) \ar[d] \ar[r] &
0 \ar[d] \\
H^p(U, \colim_i \mathcal{I}_i) \ar[r] &
H^p(U, \colim_i \mathcal{Q}_i) \ar[r] &
H^{p + 1}(U, \colim_i \mathcal{F}_i) \ar[r] &
0
}
$$
The zero at the lower right corner comes from the claim and the
zero at the upper right corner comes from the fact that the sheaves
$\mathcal{I}_i$ are injective.
The top row is exact by an application of
Algebra, Lemma \href{https://stacks.math.columbia.edu/tag/00DB}{lemma directed colimit exact (Tag 00DB)}.
Hence by the snake lemma we deduce the
result for $p + 1$.

\medskip\noindent
It remains to show that the claim is true. We will use
Lemma \href{https://stacks.math.columbia.edu/tag/03F9}{lemma cech vanish collection (Tag 03F9)}.
By the result for $p = 0$ we see that for $\mathcal{U} \in \text{Cov}$
we have
$$
\check{\mathcal{C}}^\bullet(\mathcal{U}, \colim_i \mathcal{I}_i)
=
\colim_i \check{\mathcal{C}}^\bullet(\mathcal{U}, \mathcal{I}_i)
$$
because all the $U_{j_0} \times_U \ldots \times_U U_{j_p}$
are in $\mathcal{B}$. By
Lemma \href{https://stacks.math.columbia.edu/tag/03AW}{lemma injective trivial cech (Tag 03AW)}
each of the complexes in the colimit of {\v C}ech complexes is
acyclic in degree $\geq 1$. Hence by
Algebra, Lemma \href{https://stacks.math.columbia.edu/tag/00DB}{lemma directed colimit exact (Tag 00DB)}
we see that also the {\v C}ech complex
$\check{\mathcal{C}}^\bullet(\mathcal{U}, \colim_i \mathcal{I}_i)$
is acyclic in degrees $\geq 1$. In other words we see that
$\check{H}^p(\mathcal{U}, \colim_i \mathcal{I}_i) = 0$
for all $p \geq 1$. Thus the assumptions of
Lemma \href{https://stacks.math.columbia.edu/tag/03F9}{lemma cech vanish collection (Tag 03F9)}.
are satisfied and the claim follows.
\end{proof}


\begin{lemma}
\label{lemma-colim-global}
Let $\mathcal{C}$ be a site. Let $S \subset \Ob(\Sh(\mathcal{C}))$
be a subset. Denote $*$ the final object of $\Sh(\mathcal{C})$. Assume
\begin{enumerate}
\item for some $K \in S$ the map $K \to *$ is surjective,
\item given a surjective map of sheaves $\mathcal{F} \to K$ with $K \in S$
there exists a $K' \in S$ and a map $K' \to \mathcal{F}$ such
that the composition $K' \to K$ is surjective,
\item given $K, K' \in S$ there is a surjection $K'' \to K \times K'$
with $K'' \in S$,
\item given $a, b : K \to K'$ with $K, K' \in S$ there exists a
surjection $K'' \to \text{Equalizer}(a, b)$ with $K'' \in S$, and
\item every $K \in S$ is quasi-compact
(Sites, Definition \ref{sites-definition-quasi-compact-topos}).
\end{enumerate}
Then for all $p \geq 0$ the map
$$
\colim_\lambda H^p(\mathcal{C}, \mathcal{F}_\lambda)
\longrightarrow
H^p(\mathcal{C}, \colim_\lambda \mathcal{F}_\lambda)
$$
is an isomorphism for every filtered diagram
$\Lambda \to \textit{Ab}(\mathcal{C})$, $\lambda \mapsto \mathcal{F}_\lambda$.
\end{lemma}

\begin{proof}
We will prove this by induction on $p$. The base case $p = 0$
follows from
Sites, Lemma \href{https://stacks.math.columbia.edu/tag/0GMR}{lemma directed colimits global sections (Tag 0GMR)} part (4).
We check the assumptions hold, but we urge the reader to skip this part.
Suppose $\mathcal{F} \to *$ is surjective.
Choose $K \in S$ and $K \to *$ surjective as in (1). Then
$\mathcal{F} \times K \to K$ is surjective. Choose
$K' \to \mathcal{F} \times K$ with $K' \in S$ and $K' \to K$
surjective as in (2).
Then there is a map $K' \to \mathcal{F}$ and $K' \to *$ is surjective.
Hence 
Sites, Lemma \href{https://stacks.math.columbia.edu/tag/0GMR}{lemma directed colimits global sections (Tag 0GMR)}
assumption (4)(a) is satisfied.
By Sites, Lemma \href{https://stacks.math.columbia.edu/tag/0GMP}{lemma image of quasi compact (Tag 0GMP)}, assumptions
(3) and (5) we see that $K \times K$ is quasi-compact
for all $K \in S$.
Hence 
Sites, Lemma \href{https://stacks.math.columbia.edu/tag/0GMR}{lemma directed colimits global sections (Tag 0GMR)}
assumption (4)(b) is satisfied.
This finishes the proof of the base case.

\medskip\noindent
Induction step. Assume the result holds for
$H^p$ for $p \leq p_0$ and for all topoi $\Sh(\mathcal{C})$
such that a set $S \subset \Ob(\Sh(\mathcal{C}))$ can be found
satisfying (1) -- (5). Arguing exactly as in the proof of
Lemma \ref{lemma-colim-works-over-collection}
we see that it suffices to show: given a filtered colimit
$\mathcal{I} = \colim \mathcal{I}_\lambda$ with
$\mathcal{I}_\lambda$ injective abelian sheaves, we have
$H^{p_0 + 1}(\mathcal{C}, \mathcal{I}) = 0$.
Choose $K \to *$ surjective with $K \in S$ as in (1).
Denote $K^n$ the $n$-fold self product of $K$.
Consider the spectral sequence
$$
E_1^{p, q} = H^q(K^{p + 1}, \mathcal{I}) \Rightarrow
H^{p + q}(*, \mathcal{I}) = H^{p + q}(\mathcal{C}, \mathcal{I})
$$
of Lemma \href{https://stacks.math.columbia.edu/tag/079Z}{lemma cech to cohomology sheaf sets (Tag 079Z)}. Recall that
$H^q(K^{p + 1}, \mathcal{F}) = H^q(\mathcal{C}/K^{p + 1}, j^{-1}\mathcal{F})$,
for any abelian sheaf on $\mathcal{C}$, see
Lemma \href{https://stacks.math.columbia.edu/tag/07A0}{lemma cohomology on sheaf sets (Tag 07A0)}. We have
$j^{-1}\mathcal{I} = \colim j^{-1}\mathcal{I}_\lambda$
as $j^{-1}$ commutes with colimits. The restrictions
$j^{-1}\mathcal{I}_\lambda$ are injective abelian sheaves
on $\mathcal{C}/K^{p + 1}$ by Lemma \href{https://stacks.math.columbia.edu/tag/03F3}{lemma cohomology of open (Tag 03F3)}.
Below we will show that the induction hypothesis applies
to $\mathcal{C}/K^{p + 1}$ and hence we see that
$H^q(K^{p + 1}, \mathcal{I}) = \colim H^q(K^{p + 1}, \mathcal{I}_\lambda) = 0$
for $q < p_0 + 1$ (vanishing as $\mathcal{I}_\lambda$ is injective).
It follows that
$$
H^{p_0 + 1}(\mathcal{C}, \mathcal{I}) =
H^{p_0 + 1}\left(\ldots \to
H^0(K^{p_0}, \mathcal{I}) \to
H^0(K^{p_0 + 1}, \mathcal{I}) \to
H^0(K^{p_0 + 2}, \mathcal{I}) \to \ldots\right)
$$
Again using the induction hypothesis, the complex depicted on the right
hand side is the colimit over $\Lambda$ of the complexes
$$
\ldots \to
H^0(K^{p_0}, \mathcal{I}_\lambda) \to
H^0(K^{p_0 + 1}, \mathcal{I}_\lambda) \to
H^0(K^{p_0 + 2}, \mathcal{I}_\lambda) \to \ldots
$$
These complexes are exact as $\mathcal{I}_\lambda$ is an
injective abelian sheaf (follows from the spectral sequence
for example). Since filtered colimits are exact in the category
of abelian groups we obtain the desired vanishing.

\medskip\noindent
We still have to show that the induction hypothesis applies
to the site $\mathcal{C}/K^n$ for all $n \geq 1$. Recall
that $\Sh(\mathcal{C}/K^n) = \Sh(\mathcal{C})/K^n$, see
Sites, Lemma \href{https://stacks.math.columbia.edu/tag/0791}{lemma localize topos site (Tag 0791)}.
Thus we may work in $\Sh(\mathcal{C})/K^n$.
Denote $S_n \subset \Ob(\Sh(\mathcal{C}/K^n)$ the set
of objects of the form $K' \to K^n$. We check each property in turn:
\begin{enumerate}
\item By (3) and induction there exists a surjection
$K' \to K^n$ with $K' \in S$. Then $(K' \to K^n) \to (K^n \to K^n)$
is a surjection in $\Sh(\mathcal{C})/K^n$ and $K^n \to K^n$
is the final object of $\Sh(\mathcal{C})/K^n$. Hence (1) holds for $S_n$,
\item Property (2) for $S_n$ is an immediate consequence
of (2) for $S$.
\item Let $a : K_1 \to K^n$ and $b : K_2 \to K^n$ be in $S_n$.
Then $(K_1 \to K^n) \times (K_2 \to K^n)$ is the object
$K_1 \times_{K^n} K_2 \to K^n$ of $\Sh(\mathcal{C})/K^n$.
The subsheaf $K_1 \times_{K^n} K_2 \subset K_1 \times K_2$ is the
equalizer of $a \circ \text{pr}_1$ and $b \circ \text{pr}_2$.
Write $a = (a_1, \ldots, a_n)$ and $b = (b_1, \ldots, b_n)$.
Pick $K_3 \to K_1 \times K_2$ surjective with $K_3 \in S$;
this is possibly by assumption (3) for $\mathcal{C}$.
Pick
$$
K_4
\longrightarrow
\text{Equalizer}(K_3 \to K_1 \times K_2 \xrightarrow{a_1, b_1} K)
$$
surjective with $K_4 \in S$.
This is possible by assumption (4) for $\mathcal{C}$.
Pick
$$
K_5
\longrightarrow
\text{Equalizer}(K_4 \to K_1 \times K_2 \xrightarrow{a_2, b_2} K)
$$
surjective with $K_5 \in S$.
Again this is possible. Continue in this fashion until we get
$$
K_{3 + n}
\longrightarrow
\text{Equalizer}(K_{2 + n} \to K_1 \times K_2 \xrightarrow{a_n, b_n} K)
$$
surjective with $K_{3 + n} \in S$. By construction
$K_{3 + n} \to K_1 \times_{K^n} K_2$
is surjective. Hence $(K_{3 + n} \to K^n)$ is in $S_n$ and
surjects onto the product $(K_1 \to K^n) \times (K_2 \to K^n)$.
Thus (3) holds for $S_n$.
\item Property (4) for $S_n$ is an immediate consequence of
property (4) for $S$.
\item Property (5) for $S_n$ is a consequence of
property (5) for $S$. Namely, an object $\mathcal{F} \to K^n$ of
$\Sh(\mathcal{C})/K^n$ corresponds to a quasi-compact object of
$\Sh(\mathcal{C}/K^n)$ if and only if $\mathcal{F}$ is a
quasi-compact object of $\Sh(\mathcal{C})$.
\end{enumerate}
This finishes the proof of the lemma.
\end{proof}
\end{document}
